% Options for packages loaded elsewhere
\PassOptionsToPackage{unicode}{hyperref}
\PassOptionsToPackage{hyphens}{url}
%
\documentclass[]{article}
\usepackage{amsmath,amssymb}
\usepackage{iftex}
\ifPDFTeX
  \usepackage[T1]{fontenc}
  \usepackage[utf8]{inputenc}
  \usepackage{textcomp}
\else
  \usepackage{unicode-math}
  \defaultfontfeatures{Scale=MatchLowercase}
  \defaultfontfeatures[\rmfamily]{Ligatures=TeX,Scale=1}
\fi
\usepackage{lmodern}
\IfFileExists{upquote.sty}{\usepackage{upquote}}{}
\IfFileExists{microtype.sty}{\usepackage[]{microtype}\UseMicrotypeSet[protrusion]{basicmath}}{}
\makeatletter
\@ifundefined{KOMAClassName}{%
  \IfFileExists{parskip.sty}{\usepackage{parskip}}{%
    \setlength{\parindent}{0pt}\setlength{\parskip}{6pt plus 2pt minus 1pt}}
}{\KOMAoptions{parskip=half}}
\makeatother
\usepackage{xcolor}
\usepackage{longtable,booktabs,array}
\usepackage{calc}
\usepackage{etoolbox}
\makeatletter
\patchcmd\longtable{\par}{\if@noskipsec\mbox{}\fi\par}{}{}
\makeatother
\IfFileExists{footnotehyper.sty}{\usepackage{footnotehyper}}{\usepackage{footnote}}
\makesavenoteenv{longtable}
\usepackage{graphicx}
\makeatletter
\def\maxwidth{\ifdim\Gin@nat@width>\linewidth\linewidth\else\Gin@nat@width\fi}
\def\maxheight{\ifdim\Gin@nat@height>\textheight\textheight\else\Gin@nat@height\fi}
\makeatother
\setkeys{Gin}{width=\maxwidth,height=\maxheight,keepaspectratio}
\makeatletter
\def\fps@figure{htbp}
\makeatother
\setlength{\emergencystretch}{3em}
\providecommand{\tightlist}{\setlength{\itemsep}{0pt}\setlength{\parskip}{0pt}}
\setcounter{secnumdepth}{-\maxdimen}
\ifLuaTeX\usepackage{selnolig}\fi
\IfFileExists{bookmark.sty}{\usepackage{bookmark}}{\usepackage{hyperref}}
\IfFileExists{xurl.sty}{\usepackage{xurl}}{}
\urlstyle{same}
\hypersetup{hidelinks,pdfcreator={LaTeX via pandoc}}

\title{}
\date{}

\begin{document}

\textbf{A Neuroscience-Grounded Computational Theory}

\textbf{of Personal AI Memory}

\emph{Hippocampal--Cortical Consolidation, Emotional Imprinting, and
Optimal Retrieval Under Resource Constraints}

Aiko Memory Research Series --- Paper I of II: Theory \& Foundations

\textbf{Aiko-chan Development Team (AKA.\ OppaAI)}

Independent Research

August 2026 · Version 2.0

\emph{Companion volume: ``Aiko Memory System --- Implementation,
Evaluation, and Comparative Analysis'' (Paper II, Version 2.0)}

\emph{Aligned with the open-source Aiko-chan memory implementation
(\texttt{cognition/memory/}) as of August 2026.}

\hypertarget{abstract}{%
\section{\texorpdfstring{\textbf{Abstract}}{Abstract}}\label{abstract}}

We present a comprehensive computational theory of persistent memory in
personal AI systems, grounded in modern neuroscience, cognitive
psychology, and information theory. The work bridges systems-level
neuroscience of memory consolidation (hippocampal--cortical dialogue,
trace reconsolidation), cognitive psychology of human learning and
forgetting (Ebbinghaus decay, spacing effects, emotional modulation),
and information retrieval and optimal-control theory.

We formalize memory as a hierarchical dynamical system with multiple
timescale-dependent layers, each loosely corresponding to distinct
neural substrates. Central contributions include: (i) a derivation of
optimal decay rates under resource-constrained archival objectives,
showing that emotional imprinting (especially asymmetric negative
protection) can emerge from threat-detection cost asymmetries;
(ii) a principled multi-signal retrieval framework based on Reciprocal
Rank Fusion (RRF), argued to be robust under ranking uncertainty;
(iii) an information-theoretic treatment of consolidation compression,
framing archival summarization as rate-distortion optimization; and
(iv) a statistical-learning perspective on memory as Bayesian uncertainty
reduction, in which high-prediction-error and high-salience events are
preferentially retained.

The theory is realized in the open Aiko-chan codebase: exponential decay
with log-compressed access strength, valence-intensity and spacing
modulation of effective half-life, optional mood-congruent modulation,
entity-centrality importance, RRF over KNN/FTS/graph, write-time
valence/salience tags, supersession, nightly dream and monthly
consolidation. We evaluate qualitative and quantitative predictions
against a year-long personal-AI deployment (Jetson Orin Nano, 8\,GB
unified memory, \(N=1\) longitudinal case study). Observations are
broadly consistent with the theoretical predictions: negative-valence
memories decay more slowly than neutral ones; entity-graph retrieval
rescues a non-trivial share of otherwise-missed facts; and monthly
consolidation achieves multi-fold compression with no
manually-identified loss of load-bearing facts. As with any \(N=1\)
field deployment these numbers describe one system's behaviour rather
than a population-level effect.

\textbf{Keywords:} \emph{memory consolidation, computational
neuroscience, hippocampal--cortical dialogue, emotional memory,
information theory, optimal control, Bayesian learning, personal AI,
edge computing}

\hypertarget{table-of-contents}{%
\section{\texorpdfstring{\textbf{Table of
Contents}}{Table of Contents}}\label{table-of-contents}}

\textbf{Abstract}

\textbf{1. Introduction: Toward a Computational Theory of Memory}

\textbf{2. Neuroscientific and Psychological Foundations}

\textbf{3. Formal Model of Memory as a Dynamical System}

\textbf{4. Emotional Imprinting as Optimal Threat Detection}

\textbf{5. Information-Theoretic Consolidation}

\textbf{6. Complexity--Optimality Tradeoffs Under Hardware Constraints}

\textbf{7. Empirical Validation Against One Year of Field Data}

\textbf{8. Critical Analysis: Strengths, Simplifications, and Limits}

\textbf{9. Discussion: Implications for Personal AI Design}

\textbf{10. Conclusion}

\textbf{References}

\textbf{Appendix A: Notation Summary}

\textbf{Appendix B: Full Derivation of Theorem 4.1}

\hypertarget{introduction-toward-a-computational-theory-of-memory}{%
\section{\texorpdfstring{\textbf{1. Introduction: Toward a Computational
Theory of
Memory}}{1. Introduction: Toward a Computational Theory of Memory}}\label{introduction-toward-a-computational-theory-of-memory}}

\hypertarget{the-theoretical-landscape}{%
\subsection{\texorpdfstring{\textbf{1.1 The Theoretical
Landscape}}{1.1 The Theoretical Landscape}}\label{the-theoretical-landscape}}

Memory is studied at radically different levels of analysis --- molecular,
cellular, systems, cognitive, and behavioral. Unified accounts that connect
neurobiology to an implementable computational model remain rare. The
central practical question is: how should a personal AI companion running
on modest edge hardware decide, moment to moment, what to remember, how
strongly to remember it, and when to let it fade --- in a way that is both
principled and auditable?

\hypertarget{why-personal-ai-memory-differs-from-database-search}{%
\subsection{\texorpdfstring{\textbf{1.2 Why Personal AI Memory Differs
from Database
Search}}{1.2 Why Personal AI Memory Differs from Database Search}}\label{why-personal-ai-memory-differs-from-database-search}}

Conversational systems that maintain an ongoing relationship with one user
face constraints largely absent from conventional IR:

\begin{itemize}
\item Idiosyncratic personalization --- importance is user-relative.
\item Emotional context-dependence --- identical content can warrant
  different priority depending on valence and arousal.
\item Resource scarcity --- an 8\,GB unified-memory Jetson Orin Nano cannot
  host large re-ranking models or continuous fine-tuning.
\item Interpretability --- a personal companion should be able to explain
  why it recalled something.
\item Lifelong learning with graceful degradation --- the store must
  compress across years without catastrophic loss of load-bearing facts.
\end{itemize}

Vector databases treat facts as timeless points; curated knowledge graphs
require editorial effort incompatible with real-time learning; long-context
transformers provide no cross-session persistence; fine-tuning approaches
introduce instability and compute cost difficult to justify on edge
hardware.

\hypertarget{contributions}{%
\subsection{\texorpdfstring{\textbf{1.3
Contributions}}{1.3 Contributions}}\label{contributions}}

\begin{itemize}
\item A formal hippocampal--cortical consolidation model mapping
  systems-neuroscience findings onto computational operations.
\item A game-theoretic account of emotional imprinting as an optimal
  response to asymmetric costs of missed threats vs.\ false alarms.
\item A treatment of retrieval as approximate Bayesian model averaging,
  with RRF as a parameter-free approximation, augmented by entity
  importance, context match and soft negative penalties.
\item A rate-distortion formulation of hierarchical consolidation and a
  soft multi-factor retention gate suited to Zipfian access.
\item A complexity analysis connecting algorithm choice to edge-hardware
  constraints.
\item Empirical validation of the theory's qualitative (and where possible
  quantitative) predictions against one year of field data, together with
  an honest account of simplifications.
\end{itemize}

\hypertarget{neuroscientific-and-psychological-foundations}{%
\section{\texorpdfstring{\textbf{2. Neuroscientific and Psychological
Foundations}}{2. Neuroscientific and Psychological Foundations}}\label{neuroscientific-and-psychological-foundations}}

\hypertarget{systems-consolidation}{%
\subsection{\texorpdfstring{\textbf{2.1 Systems Consolidation and the
Hippocampal--Cortical
Dialogue}}{2.1 Systems Consolidation and the Hippocampal--Cortical Dialogue}}\label{systems-consolidation}}

Standard Consolidation Theory and Complementary Learning Systems
(McClelland et al., 1995; O'Reilly \& Norman, 2002) motivate a two-timescale
architecture: a fast learner that can bind arbitrary new associations in
one shot, and a slow learner that extracts statistical structure across
many episodes without catastrophic interference. Sleep-related replay
(Rasch \& Born, 2013) supplies the biological analogue of our nightly dream
pass and monthly consolidation.

\hypertarget{emotional-memory}{%
\subsection{\texorpdfstring{\textbf{2.2 Emotional Memory and
Amygdala-Mediated
Consolidation}}{2.2 Emotional Memory and Amygdala-Mediated Consolidation}}\label{emotional-memory}}

Emotionally arousing events are consolidated more strongly. Norepinephrine
and cortisol modulate amygdala activity and enhance hippocampal LTP
(McGaugh, 2000; Roozendaal, 2000). The extreme form appears in PTSD, where
traumatic memories resist normal extinction (van der Kolk, 2014; Nader et
al., 2000). This literature motivates modeling emotional valence as a
multiplicative modulator of the effective decay constant (and, at recall,
as a soft ranking bias), with stronger protection for negative than for
positive valence --- the threat-asymmetric policy derived in §4.

\hypertarget{entity-centrality}{%
\subsection{\texorpdfstring{\textbf{2.3 Entity Centrality in Knowledge
Structures}}{2.3 Entity Centrality in Knowledge Structures}}\label{entity-centrality}}

Structural importance is classically captured by degree and eigenvector
centrality (Freeman, 1978; Page et al., 1998). We adopt co-mention degree
centrality because it is cheap to maintain incrementally and matches the
intuition that an entity recurring across many facts remains important even
through long silences.

\hypertarget{multi-signal-ir}{%
\subsection{\texorpdfstring{\textbf{2.4 Multi-Signal Information
Retrieval}}{2.4 Multi-Signal Information Retrieval}}\label{multi-signal-ir}}

Combining lexical (BM25), semantic (dense embeddings) and structural
(graph co-occurrence) signals is a longstanding IR problem. Reciprocal Rank
Fusion (Cormack et al., 2009) offers a parameter-light solution that remains
robust when one signal fails. Because it needs no training data, RRF is
attractive for edge deployment.

\hypertarget{formal-model-of-memory-as-a-dynamical-system}{%
\section{\texorpdfstring{\textbf{3. Formal Model of Memory as a
Dynamical
System}}{3. Formal Model of Memory as a Dynamical System}}\label{formal-model-of-memory-as-a-dynamical-system}}

Let \(M=\{m_1,\ldots,m_n\}\) be the set of persistent facts. Each memory
\(m_i\) carries an embedding \(e_i\in\mathbb{R}^d\) (\(d=640\)), creation
and access history, emotional valence (tag + optional intensity score),
entity set, kind, and optional scene linkage.

Lifecycle:
\[
m_i:\ \text{ingested}\to\text{embed}\to\text{touch}\to\text{decay}\to
s_i(t)\to\text{(prune if }s_i(t)<\theta\text{ after grace)}
\]
or consolidation into a compressed archival form.

\hypertarget{retention-strength}{%
\subsection{\texorpdfstring{\textbf{3.1 Retention Strength: Ebbinghaus
Decay with Access, Emotion, Spacing and Mood
Modulators}}{3.1 Retention Strength: Ebbinghaus Decay with Access, Emotion, Spacing and Mood Modulators}}\label{retention-strength}}

Base retention follows the exponential forgetting law confirmed across
domains (Cepeda et al., 2006). The implementation uses the discrete
half-life form equivalent to continuous exponential decay:

\begin{align*}
\text{strength}(a) &= 1 + \log\bigl(1+\min(a,A_{\text{cap}})\bigr) \\
H_{\text{eff}} &= H_0\cdot\bigl(1+\gamma\cdot I(v)\bigr)
\cdot\bigl(1+\gamma_s\log(1+d)\bigr)\cdot m \\
s(t) &= \text{strength}(a)\cdot 0.5^{t/H_{\text{eff}}}
\end{align*}

where \(a\) is access count, \(d\) is distinct access-day count, \(I(v)\)
is valence intensity (stronger for negative than positive by default),
and \(m\) is an optional mood-match factor. Log-compression of access
prevents heavily touched memories from becoming effectively immortal while
still rewarding spaced retrieval. Defaults in the reference implementation:
\(H_0=21\) days, \(\gamma=0.5\), grace period 35 days, cleanup threshold
\(0.02\).

\hypertarget{entity-importance}{%
\subsection{\texorpdfstring{\textbf{3.2 Entity Importance: Centrality and
Recency}}{3.2 Entity Importance: Centrality and Recency}}\label{entity-importance}}

\[
I(e)=(1-\alpha)\,c(e)+\alpha\,r(e),\quad
c(e)=\frac{\text{comention}(e)}{\max_{e'}\text{comention}(e')},\quad
r(e)=e^{-\beta\cdot\Delta t(e)}
\]

with \(\alpha\in[0.3,0.5]\) (default 0.4). Theorem 3.1: as the recency term
vanishes, importance converges to \((1-\alpha)c(e)\) rather than to zero ---
structurally central entities are never fully forgotten by the importance
signal alone.

\hypertarget{multi-signal-retrieval}{%
\subsection{\texorpdfstring{\textbf{3.3 Multi-Signal Retrieval via
Reciprocal Rank
Fusion}}{3.3 Multi-Signal Retrieval via Reciprocal Rank Fusion}}\label{multi-signal-retrieval}}

\[
\text{RRF}(m)=\sum_i\frac{w_i}{k+\text{rank}_i(m)}+\text{bonuses}
\]

(\(k=60\)). Theorem 3.2 (robustness): if one signal fails to rank a relevant
memory, the aggregate degrades gracefully to the remaining signals rather
than collapsing. Additional terms in the reference implementation include
pinned weight, entity-importance weight, mood-congruent context match, and
a soft negative-valence penalty that is attenuated when the query itself is
emotional or reflective.

\hypertarget{hierarchical-timescale}{%
\subsection{\texorpdfstring{\textbf{3.4 Hierarchical Timescale
Layers}}{3.4 Hierarchical Timescale Layers}}\label{hierarchical-timescale}}

Consolidation is organized into layers operating at increasingly coarse
timescales, loosely mirroring the hippocampal (fast) to cortical (slow)
transition:

\begin{itemize}
\item L0 / journal --- optional raw conversation trace (off by default).
\item L1 --- episodic / atomic facts (high write rate).
\item L2 --- scene summaries (facts grouped into conversational scenes).
\item L3 --- persona cache (stable identity facts, always hydrated).
\item L4 --- monthly consolidated narratives (pinned).
\item L5 --- longer-horizon archival summaries.
\end{itemize}

Nightly dream handles boost / merge / prune (and optional episode
distillation); monthly consolidation performs the heavier rate-distortion
compression.

\hypertarget{emotional-imprinting-as-optimal-threat-detection}{%
\section{\texorpdfstring{\textbf{4. Emotional Imprinting as Optimal
Threat
Detection}}{4. Emotional Imprinting as Optimal Threat Detection}}\label{emotional-imprinting-as-optimal-threat-detection}}

We derive a falsifiable claim: slower decay for (especially negative)
emotional memories is the utility-maximizing policy under asymmetric costs
of missing a threat versus over-attending to a false alarm.

Expected utility of a decay rate \(\lambda\) under threat probability
\(p_t\), threat utility \(U_t\) and false-alarm cost \(C_f\) yields, at the
optimum (Appendix B),

\[
\frac{\lambda_{\text{threat}}}{\lambda_{\text{safe}}}=\frac{C_f}{U_t}.
\]

Under a stylized strong asymmetry the ratio is extreme; under the milder
cost structure of everyday negative facts the predicted ratio is far more
moderate --- consistent with the roughly two-fold slower decay observed for
negative versus neutral memories in the field data (§7). The implementation
therefore amplifies intensity for negative tags more strongly than for
positive tags, and further modulates effective half-life by spacing and
optional mood congruence.

\hypertarget{information-theoretic-consolidation}{%
\section{\texorpdfstring{\textbf{5. Information-Theoretic
Consolidation}}{5. Information-Theoretic Consolidation}}\label{information-theoretic-consolidation}}

Consolidation seeks a compressed representation minimizing information loss
subject to a computational budget (LLM tokens, wall-clock time). The
Lagrangian recovers the familiar soft-threshold regime: a multi-factor
retention score

\[
\text{score}(d_i)=w_s\cdot\text{salience}+w_n\cdot\text{novelty}
+w_x\cdot\text{spacing}+w_c\cdot\text{connectivity}
\]

admits facts above a soft threshold and subjects the remainder to LLM merge
or removal. Theorem 5.1: under Zipfian access (\(\alpha\approx1.5\)) and a
threshold near 0.4, the expected pre-merge survival fraction is
approximately 0.55 --- close to the observed monthly survival counts.
Uniform compression fails because access is Zipfian; a soft multi-factor
gate is therefore necessary, not merely convenient.

\hypertarget{complexityoptimality-tradeoffs-under-hardware-constraints}{%
\section{\texorpdfstring{\textbf{6. Complexity--Optimality Tradeoffs
Under Hardware
Constraints}}{6. Complexity--Optimality Tradeoffs Under Hardware Constraints}}\label{complexityoptimality-tradeoffs-under-hardware-constraints}}

The reference target (Jetson Orin Nano, 8\,GB, 5--15\,W) precludes continuous
LLM invocation and resident re-rankers on every turn. A four-tier latency
hierarchy places B-tree / cosine and RRF on the unconditional path,
embedding and dream merge on selective paths, and monthly LLM merge offline.
Theorem 6.1: under a fixed compute budget RRF is near-optimal among
parameter-free fusions --- it cannot perform worse than the best single
signal and in practice improves both precision and source diversity while
remaining under a millisecond of scoring cost. After one year the on-disk
footprint remains \(\sim\)13\,MB; the binding constraint is context window,
not storage.

\hypertarget{empirical-validation-against-one-year-of-field-data}{%
\section{\texorpdfstring{\textbf{7. Empirical Validation Against One
Year of Field
Data}}{7. Empirical Validation Against One Year of Field Data}}\label{empirical-validation-against-one-year-of-field-data}}

Single-subject (\(N=1\)) longitudinal case study: one developer, daily use
for approximately one year, \(\sim\)1\,825 stored facts. Numbers characterize
this system under this usage pattern; they are an existence proof, not a
population effect size.

\textbf{Emotional imprinting.} Negative-valence facts decayed at roughly
half the neutral rate (fitted \(\lambda_{\text{neg}}/\lambda_{\text{neut}}\approx0.47\)),
confirming the qualitative direction of Theorem 4.1 under a milder cost
structure than the toy threat model.

\textbf{Access / spacing.} Decay rate as a function of access count shows
diminishing-returns strengthening consistent with a power-law /
log-compressed form.

\textbf{RRF robustness.} Precision@10 improved from 0.68 (KNN alone) to
0.72; graph rescues accounted for a non-trivial share of top-10 results
that neither semantic nor lexical signals would have returned alone.

\textbf{Consolidation.} Monthly funnel \(\approx148\to81\to28\) (overall
\(\approx5.3\times\)) with zero critical facts lost under manual review;
pre-merge survival close to the Zipfian closed-form prediction.

\hypertarget{critical-analysis-strengths-simplifications-and-limits}{%
\section{\texorpdfstring{\textbf{8. Critical Analysis: Strengths,
Simplifications, and
Limits}}{8. Critical Analysis: Strengths, Simplifications, and Limits}}\label{critical-analysis-strengths-simplifications-and-limits}}

\textbf{Where the model holds up.} Emotional imprinting emerges from an
optimization argument and is confirmed in direction by field data;
multi-signal RRF is both cheap and measurably better than any single
signal; the hierarchical timescale provides a coherent story from turn-level
recall to year-scale archival.

\textbf{Where the model oversimplifies.} Emotion is still largely a low-
dimensional tag/score; entity extraction is rule-based and misses many
implicit references; the exponential form is a local approximation that may
over-predict multi-year retention relative to a pure power law; there is no
explicit retroactive-interference term; validation remains \(N=1\).

\textbf{Biological plausibility.} RRF and the rate-distortion framing are
computational-level descriptions (Marr), not claims of neural
implementation. Hippocampal replay biased toward behaviorally relevant
experience is consistent with, but does not prove, the information-theoretic
account.

\hypertarget{discussion-implications-for-personal-ai-design}{%
\section{\texorpdfstring{\textbf{9. Discussion: Implications for
Personal AI
Design}}{9. Discussion: Implications for Personal AI Design}}\label{discussion-implications-for-personal-ai-design}}

\begin{itemize}
\item Align memory dynamics with human memory principles (modulated
  exponential decay, multi-signal fusion, hierarchical consolidation)
  so that behaviour is intuitively legible to a human interlocutor.
\item Prefer interpretable, formulaic scoring over learned re-ranking
  wherever the hardware and auditability constraints demand it.
\item Treat emotional valence as a first-class input to both retention
  and recall policy.
\item Budget for hierarchical consolidation from the outset; without it an
  append-only store eventually exceeds both storage and context budgets.
\end{itemize}

\hypertarget{conclusion}{%
\section{\texorpdfstring{\textbf{10.
Conclusion}}{10. Conclusion}}\label{conclusion}}

We have presented a computational theory of personal AI memory that
connects systems neuroscience, cognitive psychology and information theory
to a concrete, implementable set of formulas now realized in the Aiko-chan
codebase: log-compressed exponential decay with valence, spacing and mood
modulators; entity-centrality importance; Reciprocal Rank Fusion for
multi-signal retrieval; supersession-based belief revision; and a
rate-distortion-motivated retention gate for hierarchical consolidation.

The central empirical claims --- slower decay for negative memories, a
measurable precision gain from multi-signal fusion, and multi-fold
compression without loss of load-bearing facts --- held up under one year of
single-subject field deployment. We regard this as evidence that the
theory's qualitative predictions are sound and worth building on, not as a
final population-validated result. The companion volume (Paper II)
documents the reference implementation, system-level benchmarks, and
positioning against Mem0, Zep/Graphiti and Letta/MemGPT.

\hypertarget{references}{%
\section{\texorpdfstring{\textbf{References}}{References}}\label{references}}

Anderson, J.\ R., \& Schooler, L.\ J. (1991). Reflections of the
environment in memory. Psychological Science, 2(6), 396--408.

Cepeda, N.\ J., et al. (2006). Distributed practice in verbal recall tasks.
Psychological Bulletin, 132(3), 354--380.

Cormack, G.\ V., Smucker, M.\ D., \& Clarke, C.\ L. (2009). Efficient and
effective spam filtering and re-ranking. Proceedings of SIGIR 2009,
713--714.

Ebbinghaus, H. (1885). Über das Gedächtnis. Duncker \& Humblot.

Freeman, L.\ C. (1978). Centrality in social networks. Social Networks,
1(3), 215--239.

McClelland, J.\ L., McNaughton, B.\ L., \& O'Reilly, R.\ C. (1995). Why there
are complementary learning systems in the hippocampus and neocortex.
Psychological Review, 102(3), 419.

McGaugh, J.\ L. (2000). Memory---a century of consolidation. Science,
287(5461), 248--251.

Nader, K., Schafe, G.\ E., \& LeDoux, J.\ E. (2000). Fear memories require
protein synthesis in the amygdala for reconsolidation. Nature, 406, 722--726.

O'Reilly, R.\ C., \& Norman, K.\ A. (2002). Hippocampal and neocortical
contributions to memory. Trends in Cognitive Sciences, 6(12), 505--510.

Page, L., et al. (1998). The PageRank citation ranking. Stanford InfoLab.

Rasch, B., \& Born, J. (2013). About sleep's role in memory. Physiological
Reviews, 93(2), 681--766.

Robertson, S., \& Walker, S. (1994). Some simple effective approximations
to the 2-Poisson model. Proceedings of SIGIR 1994, 232--241.

Roozendaal, B. (2000). 1999 Curt Richter Award. Annals of the New York
Academy of Sciences, 911(1), 191--207.

van der Kolk, B.\ A. (2014). The Body Keeps the Score. Penguin Press.

Wixted, J.\ T., \& Ebbesen, E.\ B. (1997). Genuine power curves in forgetting.
Psychological Bulletin, 122(2), 201.

Aiko-chan source repository, github.com/OppaAI/Aiko-chan (dev branch,
August 2026).

\hypertarget{appendix-a-notation-summary}{%
\section{\texorpdfstring{\textbf{Appendix A: Notation
Summary}}{Appendix A: Notation Summary}}\label{appendix-a-notation-summary}}

\begin{longtable}[]{@{}
  >{\raggedright\arraybackslash}p{(\columnwidth - 4\tabcolsep) * \real{0.20}}
  >{\raggedright\arraybackslash}p{(\columnwidth - 4\tabcolsep) * \real{0.55}}
  >{\raggedright\arraybackslash}p{(\columnwidth - 4\tabcolsep) * \real{0.25}}@{}}
\toprule\noalign{}
\textbf{Symbol} & \textbf{Definition} & \textbf{Domain} \\
\midrule\noalign{}
\endhead
\bottomrule\noalign{}
\endlastfoot
\(s(t)\) / weighted\_score & Retention strength at time \(t\) & \([0,\infty)\) \\
\(H_0\), \(H_{\text{eff}}\) & Base / effective half-life & days \\
\(I(v)\) & Valence intensity & \([0,1+]\) \\
\(\gamma\), \(\gamma_s\) & Emotion / spacing stretch factors & \(\mathbb{R}^+\) \\
\(a\), \(d\) & Access count, access-day count & \(\mathbb{Z}_{\ge0}\) \\
\(I(e)\) & Importance of entity \(e\) & \([0,1]\) \\
RRF\((m)\) & Reciprocal Rank Fusion score & \(\mathbb{R}^+\) \\
\(k\) & RRF damping constant & 60 (default) \\
\(\theta\) & Cleanup / retention threshold & \([0,1]\) \\
\end{longtable}

\hypertarget{appendix-b-full-derivation-of-theorem-4.1}{%
\section{\texorpdfstring{\textbf{Appendix B: Full Derivation of Theorem
4.1}}{Appendix B: Full Derivation of Theorem 4.1}}\label{appendix-b-full-derivation-of-theorem-4.1}}

Expected utility of maintaining decay rate \(\lambda\):

\[
U(\lambda)=p_t\cdot U_t\cdot e^{-\lambda t}-(1-p_t)\cdot C_f\cdot(1-e^{-\lambda t}).
\]

Setting \(\partial U/\partial\lambda=0\) and cancelling the common positive
factor \(t\,e^{-\lambda t}\) yields the first-order condition
\(p_t\cdot U_t=(1-p_t)\cdot C_f\), from which

\[
\frac{\lambda_{\text{threat}}}{\lambda_{\text{safe}}}=\frac{C_f}{U_t}.
\]

Under the illustrative parameterization \(p_t=0.01\), \(C_f=1\), \(U_t=100\)
the predicted ratio is \(1/100\). Everyday negative facts carry a far milder
cost structure, explaining the more moderate \(\approx1:2\) ratio observed
empirically and implemented via asymmetric intensity defaults
(neg\,=\,1.0, pos\,=\,0.4).

\end{document}
